% =============================================================================
% SECCION: GESTION DE RIESGO
% =============================================================================

\subsection{Gesti\'on de Riesgo}
\label{sec:risk_management}

Una rentabilidad elevada solo resulta valiosa si el riesgo asumido para obtenerla se mantiene dentro de l\'imites razonables, dado que exposiciones excesivas pueden comprometer la viabilidad de una estrategia por m\'as atractivos que parezcan sus retornos. Por ello, en esta secci\'on se examina el perfil de riesgo de los cuatro modelos ganadores (LSTM+Attention, Ridge, RandomForest y CNN-LSTM), con el prop\'osito de valorar en qu\'e medida los retornos registrados guardan proporci\'on con los riesgos asumidos durante el per\'iodo de prueba.

\subsubsection{M\'etricas de Riesgo Ajustado}
\label{subsec:risk_adjusted_metrics}

Se emplean tres ratios ampliamente utilizados en gesti\'on de portafolios. El \textit{Sharpe Ratio} \citep{sharpe1966mutual} mide el exceso de retorno por unidad de volatilidad total ($\text{Sharpe} = E[R_p - R_f] / \sigma_p \cdot \sqrt{252}$). El \hyperlink{glos:sortino_ratio}{\textit{Sortino Ratio}} \citep{sortino1991} reemplaza $\sigma_p$ por la desviaci\'on est\'andar exclusivamente negativa $\sigma_d$, penalizando \'unicamente la volatilidad a la baja. El \hyperlink{glos:calmar_ratio}{\textit{Calmar Ratio}} relaciona el retorno anualizado con el m\'aximo \hyperlink{glos:drawdown}{\textit{drawdown}} hist\'orico ($\text{Calmar} = \text{CAGR} / |\text{Max DD}|$); un valor superior a 1.0 indica que el retorno anualizado supera la peor ca\'ida observada.

\begin{table}[htbp]
\centering
\caption{M\'etricas de riesgo ajustado de los modelos que superan al benchmark}
\label{tab:risk_adjusted_metrics}
\small
\begin{tabular}{@{}lrrr@{}}
\toprule
\textbf{Modelo} & \textbf{Sharpe} & \textbf{Sortino} & \textbf{Calmar} \\
\midrule
LSTM+Attention & 1.319 & 2.606 & 1.770 \\
Ridge & 0.856 & 1.250 & 0.868 \\
RandomForest & 0.651 & 1.337 & 1.153 \\
CNN-LSTM & 0.572 & 1.100 & 0.466 \\
\midrule
SPY (B\&H) & 0.657 & 0.944 & 0.571 \\
\bottomrule
\end{tabular}
\end{table}

La Tabla~\ref{tab:risk_adjusted_metrics} muestra una jerarqu\'ia consistente. LSTM+Attention registra un Sharpe de 1.319 frente al 0.657 del \textit{benchmark}, y su Sortino (2.606) es pr\'acticamente el doble de su Sharpe, lo que indica que la volatilidad proviene desproporcionadamente del lado positivo, patr\'on consistente con una estrategia que permanece en efectivo la mayor parte del tiempo y genera retornos concentrados en d\'ias de exposici\'on apalancada. LSTM+Attention y RandomForest son los \'unicos que superan el umbral de Calmar $> 1.0$, mientras que CNN-LSTM, pese a superar al \textit{benchmark} en retorno total (+122\%), presenta un Calmar (0.466) inferior al del SPY (0.571), dado que su \textit{drawdown} m\'aximo ($-35.9\%$) excede significativamente el del \'indice ($-25.4\%$).

\subsubsection{Riesgo de Cola y \textit{Drawdown}}
\label{subsec:var_cvar}

Las m\'etricas de riesgo de cola cuantifican las p\'erdidas potenciales en escenarios adversos. El VaR hist\'orico se calcula mediante percentiles emp\'iricos sobre ventanas m\'oviles de $N$ d\'ias.

\begin{equation}
\text{VaR}_{\alpha}^{(N)} = -\text{Percentil}_{\alpha}\Big(\prod_{j=0}^{N-1}(1 + R_{t-j}) - 1\Big)
\label{eq:var}
\end{equation}

El \textit{Expected Shortfall} (CVaR) mide la p\'erdida promedio condicional en las ventanas que exceden el VaR, satisfaciendo la propiedad de subaditividad que lo convierte en una medida coherente de riesgo \citep{lopezdeprado2018}. Los resultados a horizontes de 30 y 90 d\'ias (Tabla~\ref{tab:var_cvar}) muestran que LSTM+Attention presenta un VaR a 30 d\'ias ($-6.18\%$) inferior al del SPY ($-7.81\%$) y un VaR a 90 d\'ias ($-3.77\%$) que es cercano a un tercio del \textit{benchmark} ($-10.80\%$), lo que sugiere que la selectividad temporal reduce el riesgo de p\'erdidas sostenidas. En contraste, Ridge y CNN-LSTM exhiben CVaR a 90 d\'ias de $-21.88\%$ y $-23.05\%$, superando al SPY ($-13.85\%$). RandomForest exhibe el perfil m\'as conservador (VaR 90d de $-0.67\%$).

La Tabla~\ref{tab:drawdown_detail} complementa este an\'alisis. LSTM+Attention experimenta su m\'aximo \textit{drawdown} ($-20.5\%$) durante la volatilidad arancelaria de abril 2025 y se recupera en un \'unico d\'ia gracias a la posici\'on apalancada 3x durante el fuerte rebote de comienzos de abril de 2025. Ridge y CNN-LSTM registran ca\'idas del 31.8\% y 35.9\% durante el mercado bajista de 2022, con tiempos de recuperaci\'on de 56 y 566 d\'ias respectivamente. Los 566 d\'ias de CNN-LSTM reflejan las consecuencias duraderas del posicionamiento agresivo durante el per\'iodo observado. RandomForest exhibe el perfil m\'as conservador ($-15.0\%$, recuperaci\'on en 23 d\'ias).

\subsubsection{Ratios de Captura de Mercado}
\label{subsec:capture_ratios}

El \textit{Upside/Downside Capture Ratio} descompone qu\'e proporci\'on de los retornos del mercado en d\'ias alcistas y bajistas captura cada estrategia.

\begin{equation}
\text{Upside Capture} = \frac{\overline{R_p \mid R_m > 0}}{\overline{R_m \mid R_m > 0}} \times 100, \qquad \text{Downside Capture} = \frac{\overline{R_p \mid R_m < 0}}{\overline{R_m \mid R_m < 0}} \times 100
\label{eq:capture}
\end{equation}

\begin{table}[htbp]
\centering
\caption{\textit{Upside/Downside Capture} de los modelos que superan al benchmark}
\label{tab:capture_ratios}
\small
\begin{tabular}{@{}lrrr@{}}
\toprule
\textbf{Modelo} & \textbf{Upside Capture} & \textbf{Downside Capture} & \textbf{Capture Ratio} \\
\midrule
LSTM+Attention & 90.1\% & 67.2\% & 1.34$\times$ \\
Ridge & 99.5\% & 84.3\% & 1.18$\times$ \\
RandomForest & 45.5\% & 33.0\% & 1.38$\times$ \\
CNN-LSTM & 70.1\% & 61.8\% & 1.14$\times$ \\
\midrule
SPY (B\&H) & 100.0\% & 100.0\% & 1.00$\times$ \\
\bottomrule
\end{tabular}
\vspace{0.2cm}
\parbox{0.9\textwidth}{\footnotesize\textit{Nota: Calculado sobre los 1{,}301 retornos diarios evaluados del per\'iodo de prueba (710 d\'ias alcistas, 589 d\'ias bajistas, 2 d\'ias neutros). SPY B\&H = 100\% por definici\'on.}}
\end{table}

Como muestra la Tabla~\ref{tab:capture_ratios}, LSTM+Attention captura en promedio el 90.1\% de las alzas y participa del 67.2\% de las ca\'idas. Esta brecha de 22.9 puntos se acumula a trav\'es del inter\'es compuesto a lo largo del per\'iodo de prueba, lo que podr\'ia ayudar a explicar por qu\'e un \textit{Directional Accuracy} inferior al 50\% llega a asociarse con el mejor retorno del estudio. RandomForest exhibe el \textit{Capture Ratio} m\'as alto (1.38$\times$) gracias a una participaci\'on muy selectiva (45.5\% alzas, 33.0\% ca\'idas), lo que produce el \textit{drawdown} m\'as contenido del grupo pero tambi\'en el retorno m\'as modesto. CNN-LSTM presenta la menor asimetr\'ia (8 puntos de brecha, 1.14$\times$), consistente con su \textit{drawdown} m\'as profundo y su Calmar m\'as bajo, ya documentados.

El an\'alisis del \textit{Capture Ratio} complementa la paradoja del \textit{Directional Accuracy} documentada en la Secci\'on~\ref{sec:da_paradox}, dado que los modelos con mayor rentabilidad en esta muestra no parecen distinguirse por predecir correctamente la direcci\'on del mercado con mayor frecuencia, sino por registrar una asimetr\'ia entre su participaci\'on en alzas y en ca\'idas. La estrategia \textit{Long-Only} implementada incorpora un control de riesgo que se deriva del propio dise\~no del sistema de posicionamiento; cuando las predicciones no superan el umbral $q_{\text{mod}}$, la posici\'on por defecto es 100\% en efectivo, lo que operar\'ia como un \textit{stop-loss} impl\'icito, ya que reduce la exposici\'on durante per\'iodos de incertidumbre sin requerir reglas expl\'icitas predefinidas.
